% !TeX root = ../../Book3.tex
% 纲要标题：## 附录 C　非光滑度量测度空间
% 纲要位置：几何学计划纲要.md，第 3232 行。
% 正文状态：未撰写；本文件目前是框架。

% 纲要内容开始
% 
% > 与正文第29章互补；从一个实际能量和运输问题出发，不堆列合成几何缩写。
% 
% * **运输与测地结构**
%   * 耦合、Kantorovich 问题和 Wasserstein 距离
%   * 存在性的紧性方法
%   * 测度空间中的测地线
% 
% * **弱梯度与能量**
%   * 上梯度及松弛能量
%   * Sobolev 空间和 Hilbert 性
%   * 热流作为能量梯度流
% 
% * **曲率与极限**
%   * 熵凸性和曲率维数条件
%   * RCD 中附加的二次能量要求
%   * 模型和收敛稳定性
%   * 深层奇异集理论的边界
% 
% ---
% 
% 纲要内容结束

\chapter{非光滑度量测度空间}
\label{g3:app:C}
\BookChapterNumber{g3:app:C}

待撰写

\section{运输与测地结构}
\label{g3:sec:C01}
待撰写

\subsection{耦合、Kantorovich 问题和 Wasserstein 距离}
\label{g3:subsec:C01:01}
待撰写

\subsection{存在性的紧性方法}
\label{g3:subsec:C01:02}
待撰写

\subsection{测度空间中的测地线}
\label{g3:subsec:C01:03}
待撰写

\section{弱梯度与能量}
\label{g3:sec:C02}
待撰写

\subsection{上梯度及松弛能量}
\label{g3:subsec:C02:01}
待撰写

\subsection{Sobolev 空间和 Hilbert 性}
\label{g3:subsec:C02:02}
待撰写

\subsection{热流作为能量梯度流}
\label{g3:subsec:C02:03}
待撰写

\section{曲率与极限}
\label{g3:sec:C03}
待撰写

\subsection{熵凸性和曲率维数条件}
\label{g3:subsec:C03:01}
待撰写

\subsection{RCD 中附加的二次能量要求}
\label{g3:subsec:C03:02}
待撰写

\subsection{模型和收敛稳定性}
\label{g3:subsec:C03:03}
待撰写

\subsection{深层奇异集理论的边界}
\label{g3:subsec:C03:04}
待撰写
